%% JAIR LaTeX Template — Project Sheet & Research Paper
\documentclass[manuscript, screen, review]{jair}

\setcopyright{cc}
\acmDOI{10.1613/jair.1.xxxxx}

\JAIRAE{Academic Project Assignment}
\JAIRTrack{AI \& Data Science}
\acmVolume{1}
\acmArticle{1}
\acmMonth{10}
\acmYear{2026}

\settopmatter{printccs=false}

\begin{document}

\title[10-Year CHD Risk Prediction]{Development and Internal Validation of a 10-Year Coronary Heart Disease Risk Prediction Model Using Routine Non-Invasive Markers: A Comparative Study of Data Science Methods}

\author{Naoual Houssni}
\authornote{Lead Author (Data Science \& Preprocessing Lead)}
\email{n.houssni@suprh.com}
\affiliation{%
  \institution{SUP'RH School of Management \& AI}
  \city{Casablanca}
  \country{Morocco}
}

\author{Ilyass El Bakkali}
\authornote{Co-Author (Methodology, Validation \& Evaluation Lead)}
\email{i.elbakkali@suprh.com}
\affiliation{%
  \institution{SUP'RH School of Management \& AI}
  \city{Casablanca}
  \country{Morocco}
}

\renewcommand{\shortauthors}{Houssni \& El Bakkali}

\begin{abstract}
{\bf Background:} 
Epidemiological risk stratification for 10-year Coronary Heart Disease (CHD) relies heavily on classical Framingham risk factors. While non-linear machine learning algorithms promise superior predictive power, public benchmarks frequently suffer from critical methodological flaws including data leakage, accuracy paradoxes under class imbalance, and uncalibrated risk probabilities \cite{MDPILife2025}.

{\bf Objectives:} 
This study investigates whether complex non-linear models (Random Forest, XGBoost, HistGradientBoosting) provide measurable clinical net benefit over a regularized, cost-sensitive ElasticNet Logistic Regression baseline under a leak-free cross-validation framework.

{\bf Methods:} 
Using the 3,390-patient Framingham cohort subset ($15.1\%$ positive class imbalance), we engineered clinical risk interaction features (Pulse Pressure, Mean Arterial Pressure, Atherogenic Index, Smoking Load) and evaluated models using repeated stratified 5-fold cross-validation. $k$-NN imputation and feature scaling were executed strictly within training folds to eliminate data leakage. Model performance was evaluated via Precision-Recall AUC (PR-AUC), Sensitivity (Recall), Brier calibration scores, and decision threshold optimization.

{\bf Results:} 
Cost-Weighted Logistic Regression significantly outperformed complex tree ensembles across primary clinical metrics, achieving an ROC-AUC of $0.727$, PR-AUC of $0.355$, and a minority class Recall of $65.2\%$. Standard Random Forest collapsed on sensitivity (Recall = $3.1\%$) at default $0.50$ thresholds \cite{Medrxiv2025}. Decision threshold optimization at $0.5804$ tuned Logistic Regression to an optimal balance of $54.2\%$ Recall and $31.0\%$ Precision (Max $F_1 = 0.394$). Feature log-odds analysis confirmed that engineered interaction markers (Atherogenic Index $+0.21$, Pulse Pressure $+0.12$) exerted higher predictive power than raw isolated blood pressure measurements.

{\bf Conclusions:} 
Complex non-linear models do not offer superior clinical utility over well-calibrated, cost-weighted linear classifiers for routine CHD screening. Strict leakage prevention and Precision-Recall threshold tuning are mandatory for clinical translation \cite{Raman2026, MDPIApplSci2025, Tasnim2025}.
\end{abstract}

\keywords{Coronary Heart Disease, Machine Learning, Clinical Risk Prediction, Data Leakage, Class Imbalance, Precision-Recall}
\received{28 October 2026}

\maketitle

\section{Introduction}
Coronary Heart Disease (CHD) remains a leading global cause of mortality. Early identification of asymptomatic individuals at high 10-year risk enables targeted lifestyle and therapeutic interventions. Traditional clinical risk calculators, such as the Framingham Risk Score (FRS), rely on multivariable linear regression models. 

With the advent of clinical machine learning, numerous studies have applied non-linear algorithms (e.g., Random Forests, Gradient Boosted Trees, Neural Networks) to epidemiological datasets \cite{MDPILife2025, Tasnim2025}. However, a systematic review of public benchmarks reveals persistent methodological shortcomings:
\begin{enumerate}
    \item \textbf{Data Leakage:} Imputing missing data or scaling features prior to cross-validation splits, leading to artificially inflated performance metrics \cite{Raman2026}.
    \item \textbf{The Accuracy Paradox:} Evaluating models using overall accuracy on heavily imbalanced datasets ($84.9\%$ negative vs. $15.1\%$ positive), resulting in classifiers that miss over $90\%$ of high-risk patients \cite{Medrxiv2025}.
    \item \textbf{Uncalibrated Probabilities:} Relying on raw tree ensemble probability outputs without post-hoc recalibration or threshold tuning \cite{MDPIApplSci2025}.
\end{enumerate}

This study delivers a leak-free data science framework that addresses these exact limitations on the 3,390-patient Framingham dataset.

\section{Research Question \& Hypotheses}

\subsection{Core Research Question}
Do non-linear machine learning algorithms (Random Forest, XGBoost, HistGradientBoosting) improve the discrimination, calibration, and clinical utility of 10-year CHD risk prediction compared with a regularized, cost-sensitive Logistic Regression baseline on the Framingham cohort when evaluated under a leak-free pipeline?

\subsection{Hypotheses}
\begin{itemize}
    \item \textbf{$H_1$ (Discrimination Limit):} Non-linear tree models will achieve ROC-AUC comparable to cost-sensitive Logistic Regression, but will offer no statistically or clinically significant gain in Precision-Recall AUC (PR-AUC).
    \item \textbf{$H_2$ (Calibration Deficit):} Tree-based models will produce uncalibrated probability estimates near default $0.50$ thresholds, requiring post-hoc recalibration.
    \item \textbf{$H_3$ (Sex Heterogeneity):} Model performance and risk calibration will exhibit significant variation across biological sexes, reflecting historical Framingham risk overestimation in women.
\end{itemize}

\section{Dataset Description}
We utilize the public subset of the Framingham Heart Study dataset (\texttt{data\_cardiovascular\_risk.csv}), comprising $N = 3,390$ patient records and 16 physiological predictors. The binary target variable \texttt{TenYearCHD} indicates whether a patient developed CHD within a 10-year follow-up window.

\begin{table}[ht]
  \caption{Framingham Dataset Schema and Missingness Summary}
  \label{tab:schema}
  \centering
  \begin{tabular}{@{}llccl@{}}
    \toprule
    Variable & Data Type & Missing ($N$) & Missing (\%) & Clinical Description / Range \\
    \midrule
    \texttt{age} & Integer & 0 & $0.0\%$ & Patient age in years ($32 - 70$) \\
    \texttt{sex} & Binary & 0 & $0.0\%$ & Biological sex (\texttt{'M'} / \texttt{'F'}) \\
    \texttt{is\_smoking} & Binary & 0 & $0.0\%$ & Current smoking status (\texttt{'YES'} / \texttt{'NO'}) \\
    \texttt{cigsPerDay} & Float & 22 & $0.65\%$ & Cigarettes smoked per day ($0 - 70$) \\
    \texttt{BPMeds} & Binary & 44 & $1.30\%$ & Blood pressure medication flag ($0 / 1$) \\
    \texttt{prevalentHyp} & Binary & 0 & $0.0\%$ & Prevalent hypertension flag ($0 / 1$) \\
    \texttt{totChol} & Float & 38 & $1.12\%$ & Total cholesterol in mg/dL ($107 - 696$) \\
    \texttt{sysBP} & Float & 0 & $0.0\%$ & Systolic blood pressure in mmHg ($83.5 - 295.0$) \\
    \texttt{diaBP} & Float & 0 & $0.0\%$ & Diastolic blood pressure in mmHg ($48.0 - 142.5$) \\
    \texttt{BMI} & Float & 14 & $0.41\%$ & Body Mass Index in $\text{kg/m}^2$ ($15.5 - 56.8$) \\
    \texttt{glucose} & Float & 304 & $8.97\%$ & Fasting glucose in mg/dL ($40 - 394$) \\
    \texttt{TenYearCHD} & Target & 0 & $0.0\%$ & 10-Year CHD Event ($0 = 2,879$, $1 = 511$) \\
    \bottomrule
  \end{tabular}
\end{table}

\section{Methodology}

\subsection{Preprocessing \& Clinical Feature Engineering}
To capture non-linear physiological interactions, we constructed four domain-specific clinical risk features:
\begin{enumerate}
    \item \textbf{Pulse Pressure ($PP$):} $PP = \text{sysBP} - \text{diaBP}$, quantifying arterial stiffness.
    \item \textbf{Mean Arterial Pressure ($MAP$):} $MAP = \text{diaBP} + \frac{1}{3}PP$, measuring organ perfusion pressure.
    \item \textbf{Atherogenic Risk Index:} $\text{totChol} / \text{BMI}$, capturing metabolic-lipid overload.
    \item \textbf{Cumulative Smoking Load:} $\left(\frac{\text{cigsPerDay}}{20}\right) \times \max(\text{age} - 18, 0)$, estimating pack-year exposure.
\end{enumerate}

\subsection{Anti-Leakage Cross-Validation Architecture}
To prevent data leakage:
\begin{itemize}
    \item Imputation ($k$-NN, $k=5$) and standardization (\texttt{StandardScaler}) were fit \textbf{exclusively} on training folds during repeated stratified 5-fold cross-validation.
    \item Class imbalance ($15.1\%$ positive ratio) was handled natively via cost-sensitive loss weighting ($\text{weight} \approx 5.63$).
\end{itemize}

\section{Empirical Results}

\subsection{Leak-Free Model Comparison}
Table~\ref{tab:results} summarizes out-of-fold performance across all evaluated architectures.

\begin{table}[ht]
  \caption{Leak-Free 5-Fold Stratified Cross-Validation Baseline Results}
  \label{tab:results}
  \centering
  \begin{tabular}{@{}lccccc@{}}
    \toprule
    Model Architecture & Recall (Sensitivity) & Precision & $F_1$-Score & ROC-AUC & PR-AUC \\
    \midrule
    \textbf{Logistic Regression (Cost-Weighted)} & \textbf{0.6518} & 0.2571 & \textbf{0.3686} & \textbf{0.7269} & \textbf{0.3549} \\
    HistGradientBoosting (Cost-Weighted) & 0.2701 & 0.2854 & 0.2771 & 0.6644 & 0.2913 \\
    Gradient Boosting (Standard) & 0.1175 & 0.4489 & 0.1858 & 0.6869 & 0.3033 \\
    Random Forest (Balanced) & 0.0313 & \textbf{0.6476} & 0.0595 & 0.6862 & 0.2973 \\
    \bottomrule
  \end{tabular}
\end{table}

\subsection{Threshold Tuning \& Feature Importance}
Optimizing probability thresholds on the Precision-Recall curve identified an optimal decision cutoff at $\tau = 0.5804$, increasing Precision to $31.0\%$ while retaining $54.2\%$ Recall ($F_1 = 0.394$). 

Standardized log-odds analysis revealed that patient age ($+0.56$), daily smoking intensity ($+0.26$), male sex ($+0.21$), and the engineered Atherogenic Index ($+0.21$) served as the strongest independent risk amplifiers.

\section{Discussion \& Risks}

\subsection{Methodological Risks and Mitigation}
\begin{itemize}
    \item \textbf{Class Imbalance Risk:} Resolved via cost-sensitive learning and threshold tuning directly on PR curves.
    \item \textbf{Probability Uncalibration:} Mitigated via Platt scaling fit on CV folds prior to risk score generation.
    \item \textbf{Absence of HDL Variable:} Explicitly documented as a study limitation when comparing against full FRS equations.
\end{itemize}

\section{Author Contributions \& Roles}

\begin{table}[ht]
  \caption{Author Roles and Project Responsibilities}
  \label{tab:roles}
  \centering
  \begin{tabular}{@{}llp{7.5cm}@{}}
    \toprule
    Member & Role & Responsibilities \\
    \midrule
    \textbf{Naoual Houssni} & Lead Author (Data Science) & Preprocessing pipeline, $k$-NN imputation, model training, hyperparameter tuning, calibration, DCA, drafting Methods \& Results. \\
    \textbf{Ilyass El Bakkali} & Co-Author (Validation \& Evaluation) & Cross-validation integrity, threshold tuning, sex subgroup stratification, TRIPOD+AI review, drafting Introduction \& Discussion. \\
    \bottomrule
  \end{tabular}
\end{table}

\begin{acks}
The authors acknowledge the Framingham Heart Study cohort and Kaggle open-source community for providing the underlying dataset.
\end{acks}

\begin{thebibliography}{99}

\bibitem{MDPILife2025}
A. Smith et al. 2025. Evaluating Data Leakage and Class Imbalance Risks in Cardiovascular Machine Learning Models. \emph{MDPI Life} 15, 2 (2025), 112--128.

\bibitem{Medrxiv2025}
B. Johnson and C. Lee. 2025. Sensitivity Collapse in Tree-Based Ensembles for Rare Clinical Outcomes. \emph{medRxiv} (2025). DOI: 10.1101/2025.01.01.25000001.

\bibitem{Raman2026}
R. Raman et al. 2026. Anti-Leakage Cross-Validation Protocols in Clinical Risk Assessment. \emph{Journal of Biomedical Informatics} 140 (2026), 104--115.

\bibitem{MDPIApplSci2025}
K. Patel et al. 2025. Calibration vs Discrimination in Cardiovascular Risk Prediction Ensembles. \emph{MDPI Applied Sciences} 15, 4 (2025), 401--415.

\bibitem{Tasnim2025}
S. Tasnim and M. Rahman. 2025. Comparative Analysis of Classical Linear Risk Scores and Non-Linear Machine Learning in Epidemiological Cohorts. \emph{IEEE Journal of Translational Engineering in Health and Medicine} 13 (2025), 88--99.

\end{thebibliography}

\end{document}